
\section{Problem Formulation}
\label{sec:problem}

\paragraph{System Model and Constraints}
We consider two replicas A and B that must remain continuously
consistent and may reside on heterogeneous
systems. The key space is partitioned into $P$ logical units,
called \emph{segments}, by a deterministic segmenter shared by
both sides (e.g., consistent hashing or prefix mapping), so
that the same key is assigned to the same segment without
accessing engine-specific metadata.

In each reconciliation round, A and B construct and exchange a
fixed-size probe, then decode it locally to obtain diagnostic
signals that guide downstream precise verification and repair.
The probe is diagnostic rather than restorative: it does not
recover all differing items, but summarizes their scale, shape,
and likely scope.

In heterogeneous deployments, agreement on key equality and
segmentation is provided by a deployment-defined comparison
contract: a deterministic normalization layer plus a shared
segmenter. SCOUT assumes this contract and diagnoses residual
divergence with respect to it; defining the contract itself is
outside the scope of this paper. The only interface requirement
is sequential enumeration of comparable items---raw keys or
normalized identifiers that incorporate value or version
digests---with incremental hash updates and no dependence on
engine-specific internals. Because the observation budget is
much smaller than $P$, aliasing is unavoidable, and the probe
cannot guarantee useful candidate compression under all
discrepancy patterns. This limitation motivates the admission
mechanism defined below.

\paragraph{Diagnostic Primitives}
\label{sec:problem-primitives}
We now define the three output types precisely.

\textbf{Sizing (scale estimation)} estimates the global difference scale. Let $D^{+}=|A\setminus B|$ and $D^{-}=|B\setminus A|$; the total discrepancy size is $D=D^{+}+D^{-}$. Sizing always emits $\hat{D}$, regardless of gating verdict or discrepancy shape. In normal rounds, $\hat{D}$ stays near zero, with small residual values attributable to estimation noise, allowing the system to skip downstream verification; in anomalous rounds it provides a severity signal independent of the gating verdict.

\textbf{Gating (admission control)} determines whether the current round's discrepancy shape supports economical candidate compression. It emits a three-valued verdict: \texttt{notgate} admits candidate generation, \texttt{gate} diverts to a conservative fallback, and 
\texttt{unknown} (\emph{insufficient observability}) signals that the probe's fixed-budget observation lacks sufficient support for a reliable shape verdict. In this case, downstream policy combines $\hat{D}$ and other observables to determine the round's disposition.

\textbf{Triage (candidate scope)} is executed when gating
produces a \texttt{notgate} verdict. It outputs a candidate
segment set $\mathcal{C}\subseteq \{1,\ldots,P\}$ whose objective
is to minimize downstream scan scope, typically measured by
$|\mathcal{C}|$ or candidate fraction, subject to coverage of
serious discrepancies.

Together, these define the probe's interface contract: from one
digest exchange the system obtains $(g,\hat{D},\mathcal{C})$---an
admission verdict, a discrepancy-scale estimate, and a candidate
segment set produced when $g=\texttt{notgate}$. Precise
localization and repair within $\mathcal{C}$ are delegated to
downstream mechanisms such as Merkle trees or IBLTs.

\paragraph{Drift: Side Imbalance}
\label{sec:problem-drift}
The distribution of differing keys across the two sides is often
asymmetric. Under bidirectional replication, for example, both sides
may continue receiving writes while update propagation proceeds
unevenly, causing unsynchronized keys to accumulate predominantly on
one side. We refer to this systematic side imbalance as \textit{drift},
parameterized by a scalar $q\in[0,1]$ denoting side~A's share of the
total differing keys ($q=0.5$ means the differences are evenly split,
while $q\to 0$ or $1$ means they reside almost entirely on one side).
A practical probe should remain accurate across this range of
side splits rather than depending on a particular allocation of differences
between the two sides.
